%----------------------------------------------------------------------------------------
% Chapter 2

\chapter{Literature Review} % Main chapter title

\label{Chapter2} % For referencing the chapter elsewhere, use \ref{Chapter2} 
%----------------------------------------------------------------------------------------
This chapter reviews three strands of literature relevant to this thesis: first, studies on the relationship between financial systems and environmental outcomes; second, research on climate- and nature-related financial risks; and third, work on financed emissions and financial-system environmental footprints.


The relationship between finance and the environment has received increasing attention in both academic and policy discussions. One strand of the literature examines how financial systems affect environmental outcomes through the allocation of capital across firms, sectors, and technologies. Financial institutions and markets influence which economic activities receive funding and under what conditions, and therefore may shape the emissions intensity of the economy. \textcite{kim_carbon_2020}, for example, show that financial development has a non-linear relationship with CO$_2$ emissions: financial deepening may initially increase emissions by expanding economic activity, but can eventually contribute to emissions reductions as capital is reallocated towards cleaner technologies and more efficient production. They also show that financial structure matters, as more competitive and less concentrated banking systems are associated with improvements in environmental quality over the longer term. Additionally, \textcite{ecb_finance_2019} examine the role of financial structure and find that CO$_2$ emissions per capita tend to be lower in economies that rely more heavily on equity financing. A possible mechanism is that stock markets facilitate the reallocation of capital towards less polluting sectors and place greater pressure on carbon-intensive firms to adopt cleaner technologies.

The relationship also works in the opposite direction: climate change and environmental degradation can create risks for financial institutions and financial stability. The literature commonly distinguishes between physical risks, arising from the direct impacts of climate change, and transition risks, arising from the adjustment to a low-carbon economy \citep{bolton_green_2020, ngfs_call_2019}. Physical risks may damage assets, disrupt production, and reduce borrowers' repayment capacity, while transition risks may arise from carbon pricing, environmental regulation, technological change, and shifts in investor or consumer preferences \citep{ecb_financial_2021}. These risks can affect banks through credit losses and investors through asset-price adjustments, as climate-related shocks are transmitted through both microeconomic and macroeconomic channels \citep{bcbs_climate-related_2021, ecb_financial_2021}. At the aggregate level, concentrated exposures to emissions-intensive or climate-vulnerable sectors may therefore become relevant for financial stability, especially when such exposures are interconnected across institutions and markets \citep{battiston_climate_2017, ecb_financial_2021}.

More recent literature extends this climate-finance perspective to biodiversity loss and nature-related financial risks. This literature emphasizes that financial institutions are exposed not only to firms with high GHG emissions, but also to firms and sectors that depend on ecosystem services or contribute to nature degradation. \textcite{kedward_biodiversity_2023} argue that biodiversity-related financial risks are closely connected to climate-related risks, but have received comparatively less attention from central banks and supervisors. Several studies have attempted to quantify the financial-sector exposure to nature-related risks. For the Euro area, \textcite{boldrini_living_2023} find that approximately 72\% of non-financial corporations are highly dependent on at least one ecosystem service, and  almost 75\% of euro area banks' corporate loans, equivalent to nearly EUR 3.24 trillion, are granted to firms with such dependencies. Similar evidence is found at the national level: \textcite{van_toor_indebted_2020} estimate that Dutch banks, pension funds, and insurers hold EUR 510 billion in investments with high ecosystem-service dependence, corresponding to 36\% of the portfolio examined. For France, \textcite{svartzman_silent_2021} find that 42\% of the value of securities held by French financial institutions is issued by firms that are highly ecosystem-dependent, while the associated terrestrial biodiversity footprint is comparable to the loss of at least 130,000 km$^2$ of pristine nature. Recent evidence also suggests that nature-related risks may be priced in credit markets. \textcite{canipek_nature-related_2025} find that firms with higher dependence on ecosystem services face higher syndicated-loan spreads: a 1\% increase in nature dependency is associated with a 0.32\% increase in loan spreads. Taken together, these studies suggest that nature-related exposures are not marginal: they represent a sizeable share of financial institutions’ corporate lending, securities holdings, and investment portfolios.

These studies show that environmental pressures are relevant for finance through two related channels. First, financial institutions may have an impact on the environment by financing firms and sectors that generate emissions, use land and water, extract materials, or contribute to biodiversity loss. Second, financial institutions may be exposed to environmental risks when the firms they finance depend on natural capital or become vulnerable to policy, technological, legal, or market changes. This dual perspective motivates the use of environmental footprint measures that capture both the environmental impacts financed by financial institutions and the financial exposures embedded in their balance sheets and portfolios.

Within this broader context, financed-emissions metrics have become increasingly important in climate-related financial disclosure and target-setting. They are used by financial institutions to assess exposure to emissions-intensive activities and to monitor progress towards decarbonisation objectives. However, financed-emissions metrics are not always straightforward to interpret. \textcite{fraser_net-zero_2023} emphasise that reductions in financed emissions do not necessarily imply reductions in real-economy emissions. Reported financed emissions may decline because of changes in portfolio composition, ownership shares, company value, or sales, even when the underlying firms' physical emissions remain unchanged. This distinction between portfolio decarbonisation and real-economy decarbonisation is central to the interpretation of financed-emissions measures.

Despite these limitations, recent empirical evidence suggests that financed emissions are economically meaningful. \textcite{emin_financed_2025} show that credible financed-emissions targets affect bank lending behaviour: banks with such targets reduce lending to high-emission sectors and increase funding for green projects. However, firm-level emissions decline only among firms exposed to banks with more stringent targets, suggesting that financed-emissions commitments matter only when they are sufficiently credible and binding. From a financial-risk perspective, \textcite{pasiouras_financed_2025} find that banks with higher financed emissions face a higher cost of equity, indicating that investors require higher expected returns to hold shares in banks that finance more emissions-intensive activities. Together, these findings suggest that financed emissions are not only an accounting measure, but also a signal of transition-risk exposure, regulatory pressure, and market repricing.

A closely related contribution is provided by \textcite{jondeau_environmental_2026}, who develop a macroeconomic framework to measure the environmental footprint of financial systems. Instead of relying on institution-level financed-emissions disclosures, their approach combines national financial accounts with environmentally extended multi-regional input-output data. This system-level perspective complements portfolio-level approaches used by individual financial institutions and makes it possible to study financed emissions and broader environmental pressures at the level of an entire financial system. Applying this framework to Switzerland, \textcite{jondeau_environmental_2026} find that the financed emissions of the Swiss financial system amount to approximately three times Switzerland's territorial emissions. This finding highlights the importance of looking beyond domestic production when assessing the environmental footprint of financially open economies. It also shows that financial-system footprints can be substantially larger than territorial emissions when domestic financial institutions intermediate claims on foreign real-economy activity.

This thesis applies the framework of \textcite{jondeau_environmental_2026} to two major financial systems: the United States and the Euro area. While previous studies have examined the relationship between financial structure and emissions, the interpretation of financed-emissions metrics, and the financial materiality of financed emissions, less attention has been paid to the environmental pressures financed by entire financial systems. By comparing a more market-based system with a more bank-based one, this thesis contributes to understanding how the structure of financial intermediation is associated with different patterns of financed environmental pressure. It also extends the analysis beyond GHG emissions by considering broader environmental pressures, including land use, water use, and material extraction.











